\documentclass[a4paper,12pt, oneside]{article}   	
	
\usepackage[utf8x]{inputenc}
\usepackage[english,greek]{babel}
\usepackage[margin=1.2in]{geometry}
\usepackage[parfill]{parskip}    	
\usepackage{microtype}
\usepackage{graphicx}
\usepackage{wrapfig}
\usepackage{enumitem}
\usepackage{fancyhdr}
\usepackage{amsmath,amsfonts,amssymb}
\usepackage{cite}
\usepackage{subcaption}
\usepackage{float}
\usepackage{titlesec}
\usepackage{index}
\makeindex

%SetFonts
\usepackage{lmodern}
%SetFonts
\usepackage[breaklinks=true, unicode=true]{hyperref}

\pagestyle{fancy}
\fancyhf{}
\renewcommand{\headrulewidth}{0pt}
\fancyhead[RE]{}
\fancyhead[RO]{}
\fancyfoot[RE, RO]{\thepage}


\titleformat{\section}
  {\normalfont\Large\bfseries}{\thesection}{1em}{}[{\titlerule[0.8pt]}]

\begin{document}

\begin{titlepage}
\begin{center}

\includegraphics[width=0.20\textwidth]{./img/NTUAlogo.jpg}~\\[0.1cm]
\textsc{ ΕΘΝΙΚΟ ΜΕΤΣΟΒΙΟ ΠΟΛΥΤΕΧΝΕΙΟ}\\[0.2cm]  
\textsc{ ΣΧΟΛΗ ΗΛΕΚΤΡΟΛΟΓΩΝ ΜΗΧΑΝΙΚΩΝ ΚΑΙ ΜΗΧΑΝΙΚΩΝ ΥΠΟΛΟΓΙΣΤΩΝ}\\[3cm] 


\textbf{\LARGE \textlatin{Stakeholders Requirements Specification}}\\[0.01cm]
\text{Σύμφωνα με το πρότυπο \textlatin{ISO/IEC/IEEE 29148:2011}}\\[2cm]

\textbf{\Large Εμπλεκόμενο μέρος: Υπ. Ενέργειας}\\[2cm]

% Author and supervisor 
\textbf{Ομάδα:  \textlatin{WeAreBack}}\\
	Σταύρος Σταύρου \\
	Λένος Τσοκκής \\ 
	Κωνσταντίνα Φουντουραδάκη
\vfill

\end{center} 
\end{titlepage}

\thispagestyle{empty}
\pagenumbering{gobble}% Remove page numbers (and reset to 1)
\newpage

\thispagestyle{empty}
\pagenumbering{arabic}
\setcounter{page}{2}
\tableofcontents
\newpage


\section{Εισαγωγή}
\subsection{Ταυτότητα - Επιχειρησιακοί στόχοι}
Στόχος είναι η δημιουργία ενός συστήματος λογισμικού που εξασφαλίζει τη διαφάνεια στη λειτουργία της ευρωπαϊκής αγοράς ηλεκτρικής ενέργειας. Δίνει στα εμπλεκόμενα μέρη δυνατότητες διάθεσης, οπτικοποίησης και ανάλυσης \href{https://transparency.entsoe.eu}{ανοικτών δεδομένων.}

Μέσω του συστήματος αυτoύ το Υπ. Ενέργειας μπορεί να παρακολουθεί και να αναλύει τις αγορές ενέργειας και τα διατομεακά ζητήματα και να εντοπίζει πιθανές απειλές για την τακτική λειτουργία του ενεργειακού συστήματος. Επιπλέον, παρακολουθεί θέματα που σχετίζονται με την ενέργεια όπως η διασφάλιση της ασφάλειας εφοδιασμού της της χώρας του σε καύσιμα. Η εφαρμογή έχει ως στόχο να συμβάλει στην αποτελεσματική ενεργειακή πολιτική της χώρας και στη μεγαλύτερη διαφάνεια των ενεργειακών αγορών. Αποστολή της είναι η συγκέντρωση πληροφοριών και η ανάλυση των αγορών ενέργειας με σκοπό να βοηθηθούν οι υπεύθυνοι χάραξης πολιτικής να αξιολογήσουν, να προτείνουν ή να προσαρμόσουν τις πολιτικές της χώρας. \\Τέλος η εφαρμογή στοχεύει στην στήριξη της ατζέντας πολιτικής της χώρας και βοηθά το Υπουργείο να προβλέψει προβλήματα και να υποβάλει πιθανές λύσεις.


\subsection{Περίγραμμα επιχειρησιακών λειτουργιών}
Το Υπ. Ενέργειας, αφού συνδεθεί στο σύστημα μπορεί να πραγματοποιήσει αναζήτηση για τιμές κατανάλωσης ενέργειας πραγματικές ή εκτιμώμενες, να τις συγκρίνει καθώς και να αναζητήσει τις τιμές παραγωγής ενέργειας. Για αυτές τις λειτουργίες μπορεί να εφαρμόσει κάποια φίλτρα στα δεδομένα που θα του επιστραφούν, όπως περιοχή, χρονική ανάλυση, χρονικό διάστημα και τύπο ενέργειας, αλλά και να  επιλέξει το μορφότυπο τους.

\newpage
\section{Αναφορές-πηγές πληροφορίας}
\begingroup
%\renewcommand{\section}[2]{}
\renewcommand{\section}{\subsection}
\begin{thebibliography} {books}
\latintext
\bibitem{ISO2011} ISO/IEC 29148:2011 (IEEE Std 29148-2011), Systems and software engineering - Life cycle processes - Requirements engineering
\bibitem{ISO25060} ISO/IEC TR 25060:2010  Systems and software engineering - Systems and software product Quality Requirements and Evaluation - Common Industry Format for usability: General framework for usability-related information
\end{thebibliography}
\endgroup
\subsection*{Πηγές πληροφορίας}
\begingroup
\latintext
\sloppy
\begin{enumerate}
\begin{sloppypar}
   \item \href{https://transparency.entsoe.eu/content/static\_content/Static\%20content/knowledge\%20base/SFTP-Transparency\_Docs.html}{SFTP-Transparency Docs}
   \item \href{https://ec.europa.eu/energy/sites/ener/files/documents/moe_and_emos_description.pdf}{MOE and EMOS description}
   \end{sloppypar}
\end{enumerate}
\endgroup
\greektext


\section{Διαχειριστικές απαιτήσεις επιχειρησιακού περιβάλλοντος} 
\subsection{Επιχειρησιακό μοντέλο}
Η εφαρμογή θα λειτουργήσει και θα γίνει διαδεδομένη λόγω του ότι:
\begin{itemize}
  \item Η πλατφόρμα είναι αξιόπιστη καθώς εξασφαλίζει την ασφάλεια δεδομένων του εκάστοτε χρήστη μέσω κρυπτογραφίας.
  \item Το Υπουργείο έχει πρόσβαση σε μια μεγάλη βάση δεδομένων κατανάλωσης ενέργειας, που διαρκώς ενημερώνεται.
  \item Παρέχει διαφάνεια στην λειτουργία της αγοράς ηλεκτρικής ενέργειας.
    \item Αποτελεί για το Υπουργείο μέσον εντοπισμού απειλών προς το ενεργειακό σύστημα.
\end{itemize}


\subsection{Περιβάλλον διαχείρισης πληροφοριών}
Η σημερινή εικόνα για το περιβάλλον διαχείρισης πληροφοριών  στον τομέα της ηλεκτρικής ενέργειας περιλαμβάνει τους ακόλουθους οργανισμούς και πλατφόρμες. Η Ευρωπαϊκή ένωση παρέχει το σύστημα παρατήρησης \textlatin{EMOS} και ανά δύο χρόνια \href{https://ec.europa.eu/energy/en/data-analysis/market-analysis#gas-and-electricity-market-reports}{αναφορές} για την αγορά (τόσο τη βιομηχανική όσο και τα νοικοκυριά) ηλεκτρικής ενέργειας  που αναλύουν την εξέλιξη τιμών και όγκων κατανάλωσης καθώς και των αλληλεπιδράσεων μεταξύ των κρατών μελών της. Εταιρίες όπως η \href{https://www.enerdata.net}{\textlatin{Enerdata}}, \href{https://ec.europa.eu/eurostat/web/energy}{\textlatin{eurostat}}, \href{http://www.vaasaett.com}{\textlatin{vaasa ETT}} είναι μερικές από τις εταιρίες που προσφέρουν ανάλυση της αγοράς ενέργειας. 
\newpage


\section{Λειτουργικές απαιτήσεις επιχειρησιακού περιβάλλοντος}
\subsection{Επιχειρησιακές διαδικασίες}
Οι διαδικασίες της εφαρμογής που απαιτεί ένα Υπουργείο Ενέργειας περιγράφονται ακολούθως.
\renewcommand{\theenumi}{\Roman{enumi}}
\renewcommand{\theenumii}{\arabic{enumii}}
\renewcommand{\theenumiii}{\alph{enumiii}}

\begin{enumerate}
   \item Πρόσβαση - Απεικόνιση δεδομένων
      \begin{enumerate}
     \item ανά περιοχή \textlatin{(Scope)}
     \item με συγκεκριμένη χρονική ανάλυση \textlatin{(Resolution)}
     \item για δοσμένη χρονική στιγμή \textlatin{(date)}
     \begin{enumerate}
  	\item συγκεκριμένη ημερομηνία
   	\item συγκεκριμένο μήνα και έτος	   	
	\item συγκεκριμένο έτος
     \end{enumerate}
     \item με μορφότυπο \textlatin{format}
     \begin{enumerate}
  	\item \textlatin{JSON}
   	\item \textlatin{.csv} αρχείο
     \end{enumerate}
   \end{enumerate}
   \item Σύγκριση δεδομένων
       \begin{enumerate}
     \item ανά περιοχή \textlatin{(Scope)}
     \item με συγκεκριμένη χρονική ανάλυση \textlatin{(Resolution)}
     \item για δοσμένη χρονική στιγμή \textlatin{(date)}
     \begin{enumerate}
  	\item συγκεκριμένη ημερομηνία
   	\item συγκεκριμένο μήνα και έτος	   	
	\item συγκεκριμένο έτος
     \end{enumerate}
     \item με μορφότυπο \textlatin{format}
     \begin{enumerate}
  	\item \textlatin{JSON}
   	\item \textlatin{.csv} αρχείο
     \end{enumerate}
   \end{enumerate}

   \item Σύνδεση στο σύστημα
    \begin{enumerate}
     \item σύνδεση με όσο το δυνατό λιγότερα πεδία προς συμπλήρωση, όπου \\ απαραίτητα είναι τα \textlatin{username} και  \textlatin{password}
     \item ανάκτηση κωδικού πρόσβασης
    \end{enumerate}
  \item Αποσύνδεση από το σύστημα
  \item Έλεγχος συνδεσιμότητας μεταξύ του χρήστη και της βάσης δεδομένων \textlatin{(HealthCheck)}
\end{enumerate}


\subsection{Δείκτες ποιότητας}
Δείκτες ποιότητας που σηματοδοτούν τη σωστή λειτουργία της εφαρμογής είναι ο μέσος αριθμός νέων χρηστών ανά ημέρα καθώς και το μεγάλο πλήθος εγγεγραμμένων χρηστών. Επιπλέον πολύ βασική ένδειξη ποιότητας είναι  η συχνότητα επισκέψεων των χρηστών. Ακόμη, δείκτης ποιότητας της πλατφόρμας είναι το πλήθος των δεδομένων με το οποίο ενημερώνεται αλλά και το εύρος περιοχών που καλύπτει. Τέλος σημαντική ένδειξη ποιότητας συνιστά η χρήση της εφαρμογής και από άλλους διεθνείς και εθνικούς εταίρους, φορείς της αγοράς και ακαδημαϊκούς κύκλους.
\renewcommand{\theenumi}{\arabic{enumi}}
\renewcommand{\theenumii}{\alph{enumii}}
\begin{enumerate}
   \item Ασφάλεια 
   \begin{enumerate}
     \item Τα στοιχεία των εγγεγραμμένων χρηστών είναι κρυπτογραφημένα.
     \item Χρήση πρωτοκόλλου \textlatin{HTTPS} για ασφαλή δικτυακή σύνδεση στην οποία τα δεδομένα ανταλλάσονται κρυπτογραφημένα.
   \end{enumerate}
   \item Περιεχόμενο της πλατφόρμας
   \begin{enumerate}
     \item Έγκαιρο, αφού καταχωρούνται καθημερινά στην πλατφόρμα νέα δεδομένα.
     \item Έγκυρα δεδομένα, τα οποία χρησιμοποιούνται και από τα κράτη μέλη της Ευρωπαϊκής Ένωσης.
     \item Ευρείας χωρικής κάλυψης
     \item Δεδομένα μονίμως καταχωρημένα, καθώς δεν διαγράφονται όταν απαρχαιωθούν.
   \end{enumerate}
\end{enumerate}


\section{Έκθεση απαιτήσεων χρηστών}
Το Υπουργείο Ενέργειας απαιτεί από την εφαρμογή να καλύπτει όλα τα είδη ενέργειας,  να συγκεντρώνει δεδομένα όπως τιμές καυσίμων, ισοζύγιο εμπορίου - ανά ενεργειακό τύπο και ανά χώρα, φόρους επί της ενέργειας κτλ.  Ακόμη επιθυμεί εύκολη χρήση της πλατφόρμας και στενή παρακολούθηση των τιμών της ενέργειας ανά τύπο μέσω διαγραμμάτων όπως και οπτικοποίηση των υπόλοιπων δεδομένων. Τέλος επιθυμεί την ακεραιότητα του συστήματος, την ασφάλεια πρόσβασης και την προστασία δεδομένων.

\newpage
\section{Αρχές του προτεινόμενου συστήματος}
Το Υπουργείο Ενέργειας επιθυμεί διεπαφή εύκολη στη χρήση,  με αξιόπιστη ομαλή λειτουργία, ταχύτατη ανάκτηση των δεδομένων και εύμορφη οπτικοποίηση των αποτελεσμάτων. Ακόμη, απαραίτητη αρχή συνιστά η προστασία από υποκλοπή των στοιχείων τους με χρήση κρυπτογραφίας. Τέλος, τα έγγραφα τεκμηρίωσης της εφαρμογής να είναι στα ελληνικά.


\section{Περιορισμοί στο πλαίσιο του έργου}
Περιορισμοί στο πλαίσιο του έργου αποτελούν:
\begin{itemize}
  \item Δυσκολία κατανόησης της χρήσης του \textlatin{CLI}
  \item Δυσκολία κατανόησης της χρήσης του \textlatin{SSH} για σύνδεση στο \textlatin{server}
  \item Απαιτείται σύνδεση στο διαδίκτυο για πρόσβαση σε τιμές που δεν είχαν ήδη ανακτηθεί ως \textlatin{.csv} αρχείο
  \item Η γλώσσα των δεδομένων και των διεπαφών είναι η αγγλική
    \item Τα δεδομένα της εφαρμογής αφορούν μόνο την Ευρώπη
    \item Η συλλογή δεδομένων δεν έχει επικεντρωθεί στην απόλυτη κάλυψη μιας χώρας και  δε αυτής του αντίστοιχου υπουργείου 
\end{itemize}
\newpage

\section{Παράρτημα: ακρωνύμια και συντομογραφίες}
\textlatin{\textbf{EMOS} Environmental Management Overview Strategy\\
\textbf{REST API}  RESTful Application Programming Interface\\
\textbf{CLI} Command Line Interface\\
\textbf{JSON} JavaScript Object Notation\\
\textbf{HTTPS} Hypertext Transfer Protocol Secure\\
\textbf{SSH} Secure Shell \\
\textbf{GDPR} General Data Protection Regulation
}

\end{document}  